\documentclass[10pt,a4paper]{amsart}
\usepackage[margin=19mm]{geometry}
\usepackage{amsmath,amssymb,mathtools}
\usepackage[scr=rsfs,cal=euler]{mathalfa}
\usepackage{xfrac}
\usepackage[unicode,hidelinks,bookmarks=false]{hyperref}
\newcommand{\ie}{i.e.\ }
\newcommand{\cf}{cf.\ }
\newcommand{\resp}{respectively\ }
\newcommand{\Subsection}{\subsection}
\providecommand{\nobreakdash}{\nobreak\hbox{-}}
\setlength{\emergencystretch}{2em}
\multlinegap0pt
\usepackage{amssymb}
\usepackage{xcolor}
\newtheorem{lemma}{Lemma}
\def\S{\Bbb S}
\def\k{\kappa}
\def\l{\lambda}
\def\p{\partial}
\title{An isoperimetric inequality for Neumann eigenvalues with radial log-concave measures}
\author{Kui Wang, Xiaomei Sun, Anqiang Zhu}
\date{Source: arXiv:2608.08748v1 (CC BY 4.0)}
\begin{document}
\maketitle

\noindent\emph{Source and licence.} An isoperimetric inequality for Neumann eigenvalues with radial log-concave measures. Version arXiv:2608.08748v1 (CC BY 4.0), distributed by arXiv under the Creative Commons Attribution 4.0 International licence, http://creativecommons.org/licenses/by/4.0/, as recorded in the arXiv licence metadata for this exact version and shown on the abstract page https://arxiv.org/abs/2608.08748v1. Full licence notice: \texttt{licenses/arxiv-2608.08748-CC-BY-4.0.txt}.\par\smallskip

\noindent\textit{Editorial scope.} Counted unit: the article's lemma lm3.1 (source0.tex lines 353--366). The article's complete original proof of it is delivered with it (source0.tex lines 367--453, one \texttt{proof} environment of 87 lines). No mathematics is written, repaired or re-derived: the statement and the proof are the article's own lines, in the article's own notation, and no retained line is substituted or re-flowed.  Retained uncounted as the source-local context the proof needs: lines 345--350.  Nothing was omitted from any retained span, so the package carries no omission record.  No retained line contains a citation, so the wrapper carries no bibliography and no bibliography entry.  No retained line contains a figure environment, an image-inclusion command or drawing code, so no image is delivered and the package declares no asset dependency.  Editorial formatting only, in addition: an AMS \texttt{amsart} preamble that restates the article's own declarations and macros used by the retained lines, namely the environments \texttt{lemma} on separate counters, together with the standard packages those lines need.  The retained environments print in this document as Lemma 1 in order of appearance, whereas the article prints the same items with the numbering of its own full document.  The complete native source of the article as distributed in the arXiv e-print for this exact version is bundled unchanged as \texttt{sources/207185/source0.tex}.
\begin{align}\label{3.1}
   \begin{cases}
       -\Delta u+\nabla u \cdot \nabla \phi=\mu u, \quad& \text{in $B_R$,}\\
       \frac{\p u}{\p \nu}=0, \quad& \text{on $\p B_R$,}
   \end{cases} 
\end{align}

\begin{lemma}\label{lm3.1}
Let $B_R\subset M_\k^n$ be the origin-centered ball of radius $R$ with $R<\pi$ if $\k=1$. Suppose $\phi$ is  radial and convex. Then the first nonzero eigenvalue  of \eqref{3.1} has muliplicity $n$, and all the corresponding eigenfunctions are of the form
$$
u(x)=T(r)\psi_i(\theta), \quad i=1,2, \cdots, n,
$$
where $(r, \theta)$ are polar coordinates, $\psi_i(\theta)$ are the linear coordinate functions restricted to $\mathbb{S}^{n-1}$, and 
$T(r)$ is the first eigenfunction of the boundary value problem
\begin{align}\label{3.2}
    \begin{cases}
T'' + \left((n-1)\frac{C_\kappa}{S_\kappa} - \phi'\right)T' + \big(\lambda - (n-1)S_\kappa^{-2}\big)T = 0,\quad r\in (0, R),\\
T(0) = 0,\quad T'(R) = 0. 
\end{cases}
\end{align}
\end{lemma}

\begin{proof}
Since the  metric on $B_R$ admits  the polar decomposition $$g_\k=dr^2+S_\k(r)^2 g_{\S^{n-1}},$$  and $\phi$ is radial symmetric, we may apply the separation of variables
technique to find solutions of \eqref{3.1}. 
Let $u(x)=T(r)\psi(\theta)$ be an eigenfunction of \eqref{3.1}. A   direct computation yields
\begin{align}\label{3.3}
   -\frac{T'' + \left[(n-1)\frac{C_\kappa}{S_\kappa} - \phi'\right]T'}{T} -\frac 1 {S_\k^2}\frac{\Delta_{\S^{n-1}}\psi}{\psi}=\mu.
\end{align}
where $\Delta_{\S^{n-1}}$ denotes the  Laplace-Beltrami operator  on $\S^{n-1}$. Thus $\psi$ is an eigenfunction of $\Delta_{\S^{n-1}}$.
Recall that  the eigenvalues of $\Delta_{\S^{n-1}}$ are $k(k+n-2)$  for $k=0, 1, 2, \cdots$. Consequently, from  \eqref{3.3} we obtain 
\begin{align*} 
T'' + \left[(n-1)\frac{C_\kappa}{S_\kappa} - \phi'\right]T' + \left(\mu - \frac{k(k+n-2)}{S_\kappa^2}\right)T = 0.
\end{align*}
It follows that the first nonzero eigenvalue $\mu_1(B_R)$ of  \eqref{3.1}  is either the second eigenvalue $\tau_2$ of
 \begin{align}
     \begin{cases}\label{3.4}
T'' + \left[(n-1)\frac{C_\kappa}{S_\kappa} - \phi'\right]T'+ \tau T = 0, & \text{in } (0,R),\\
T'(0) = 0,\quad T'(R)  = 0,
\end{cases}
 \end{align}
or the first eigenvalue $\l_1$ of 
\begin{align}\label{3.5}
    \begin{cases}
T'' + \left[(n-1)\frac{C_\kappa}{S_\kappa} - \phi'\right]T'+ (\l-\frac{n-1}{S_\k^2})T=0, & \text{in } (0,R),\\
T(0) = 0,\quad T'(R) = 0.
\end{cases}
\end{align}
We now show $\l_1<\tau_2$.

 Let $g$ and $f$ be eigenfunctions associated to $\l_1$ and $\tau_2$, respectively.
By the Sturm-Liouville theorem,  $g$ has constant sign on on $(0,R)$ and we assume without loss of generality that $g(r)>0$ for $r\in (0,R)$. The function $f(r)$ has two nodal sets: there exists $r_{0}\in (0,R)$ such that $f(r)>0,r\in (0,r_{0})$ and $f(r)<0,r\in (r_{0},R)$.
Set
$$
p(r)=S_{\k}(r)^{n-1}e^{-\phi(r)}.
$$
Then  \eqref{3.4} can be rewritten as 
\begin{align*}
\begin{cases}
\left(p(r)f'(r)\right)' + \tau_2 \,p(r)\,f(r) = 0, \quad r\in (0, R),\\
f'(0) = f'(R) = 0,
\end{cases}
\end{align*}
Integrating over $(0,R)$ gives
\begin{align}\label{3.6}
    \int_{0}^{R}\tau p(r)f(r)dr=-p(r)f'(r)\Big|_{0}^{R}=0.
\end{align}
Since $f(r)<0$ for $r\in (r_{0},R)$, we deduce from \eqref{3.6} that
\begin{align*}
    0<\int_{0}^{r}\tau_2 p(r)f(r)dr=-p(r)f'(r), \quad r\in (0, R).
\end{align*}
 Hence $f'(r)<0 $ on $ r\in (0,R)$.
Let $h(r)=f'(r)<0$. Differentiating \eqref{3.4}, we derive
\begin{align}\label{3.7}
   h'' + h'\left(\dfrac{(n-1)C_{\kappa}}{S_{\kappa}} - \phi'\right) + \tau_2 h+(-\frac{n-1}{S_{\k}^{2}}-\phi'')h = 0, \quad \text{in } (0,R).
\end{align}
Since $\phi(r)$ is convex, we have $-\phi''(r)h(r)\geq 0$.  Thus  \eqref{3.7} implies the differential inequality
\begin{align}\label{3.8}
    h'' + h'\left(\dfrac{(n-1)C_{\kappa}}{S_{\kappa}} - \phi'\right) + \tau_2 h-\frac{n-1}{S_{\k}^{2}}h \leq  0, \quad \text{in } (0,R),
\end{align}
with boundary conditions  $h(0)=0$ and $h(R)=0$.

Comparing \eqref{3.8} with  \eqref{3.5}, we obtain 
\begin{align}\label{3.9}
    &\int_{0}^{R}\left(h''(r)g(r)-h(r)g''(r)\right)p(r)dr\nonumber\\
    +&\int_{0}^{R}\left(\dfrac{(n-1)C_{\kappa}(r)}{S_{\kappa}(r)} - \phi'(r)\right) \left(h'(r)g(r)-h(r)g'(r)\right)p(r)dr\nonumber\\
    +&\int_{0}^{R}(\tau_2-\l_1)h(r) g(r)p(r)dr\nonumber\\
    \leq& 0.
\end{align}
Integration by parts gives
\begin{align}\label{3.10}
    &\int_{0}^{R}(h''(r)g(r)-h(r)g''(r))p(r)dr\nonumber\\
    =&\int_{0}^{R}(h'(r)g(r)-h(r)g'(r))'p(r)dr\nonumber\\
    =&\int_{0}^{R}((h'(r)g(r)-h(r)g'(r))p(r))'dr-\int_{0}^{R}(h'(r)g(r)-h(r)g'(r))p'(r)dr.
\end{align}
Substituting  \eqref{3.10} into \eqref{3.9}, we have 
\begin{align} \label{3.11}
    &(\tau_2-\l_1)\int_{0}^{R}h(r)g(r)p(r)dr\nonumber\\
    \leq& -(h'(r)g(r)-h(r)g'(r))p(r)\Big|_{0}^{R}\nonumber\\
    =&-(h'(R)g(R)-h(R)g'(R))p(R)\nonumber\\
    =&-h'(R)g(R)p(R).
\end{align}
Since $h(r)\leq 0$ and $h(R)=0$, we have $h'(R)>0$.  Hence  $-h'(R)g(R)p(R)< 0$.
Because $g(r)> 0$ and $h(r)< 0$ on $r\in (0,R)$, it follows from \eqref{3.11} that
\begin{align*}
    \l_1< \tau_2.
\end{align*}
Thus $\mu_1(B_R)=\l_1$. Recalling that the first nonzero eigenvalue of $\Delta_{\S^{n-1}}$ has muliplicity $n$ with corresponding eigenfucntions  $\psi_i(\theta)$ given by  the linear coordinate functions restricted to $\S^{n-1}$, we conclude that the first nonzero eigenvalue of \eqref{3.1} has muliplicity $n$ and the corresponding eigenfucntions are given by $u_i(r, \theta)=T_1(r)\psi_i(\theta)$, where $T_1(r)$ is the first eigenfunction of \eqref{3.2}.
\end{proof}

\end{document}
